Table~\ref{tab:main} gives the main results, Fig.~\ref{fig:tradeoff} the trade-off between transport time and spill rate as the tuned scalar is varied, and Fig.~\ref{fig:traj} one example episode per system.

\begin{table*}[t]\centering
\caption{Main results at each system's tuned setting (mean $\pm$ std over 5 seeds, 50 test episodes per seed). Columns: mean transport time [s], spill rate, mean peak wall elevation [freeboards], 20-step prediction RMSE [freeboards]. Tasks: nominal liquid, nominal liquid with a push, held-out liquids, held-out liquids with a push, single-linear-mode ground truth. Bold and underline mark the best and second-best value of each column separately; a low transport time is only meaningful together with the spill rate.}\label{tab:main}
\resizebox{0.92\textwidth}{!}{\begin{tabular}{llccccc}
\toprule
Method & Task & transport\_time\_s $\downarrow$ & spill\_rate $\downarrow$ & peak\_disp $\downarrow$ & pred\_rmse\_20 $\downarrow$ & $n$ \\
\midrule
Pendulum MPC & nominal & 2.147 $\pm$ 0.012 & 0.000 $\pm$ 0.000 & \textbf{0.799 $\pm$ 0.010} & 0.135 $\pm$ 0.005 & 5 \\
ZV input shaping & nominal & \textbf{1.890 $\pm$ 0.000} & 0.000 $\pm$ 0.000 & 0.975 $\pm$ 0.007 & -- & 5 \\
Learned MPC (no context) & nominal & 1.987 $\pm$ 0.018 & \underline{0.000 $\pm$ 0.000} & 0.904 $\pm$ 0.022 & \underline{0.092 $\pm$ 0.003} & 5 \\
Spring-mass MPC & nominal & 2.041 $\pm$ 0.014 & 0.000 $\pm$ 0.000 & 0.925 $\pm$ 0.011 & 0.136 $\pm$ 0.005 & 5 \\
Two-mode spring-mass MPC & nominal & \underline{1.965 $\pm$ 0.007} & 0.000 $\pm$ 0.000 & 0.904 $\pm$ 0.007 & 0.094 $\pm$ 0.007 & 5 \\
\textbf{Learned-Ctx MPC} & nominal & 1.993 $\pm$ 0.010 & \textbf{0.000 $\pm$ 0.000} & \underline{0.887 $\pm$ 0.013} & \textbf{0.063 $\pm$ 0.004} & 5 \\
\midrule
Pendulum MPC & heldout & 2.193 $\pm$ 0.008 & \textbf{0.000 $\pm$ 0.000} & \textbf{0.800 $\pm$ 0.014} & 0.473 $\pm$ 0.015 & 5 \\
ZV input shaping & heldout & \textbf{1.890 $\pm$ 0.000} & 1.000 $\pm$ 0.000 & 1.253 $\pm$ 0.021 & -- & 5 \\
Learned MPC (no context) & heldout & 2.043 $\pm$ 0.024 & 0.100 $\pm$ 0.071 & 0.902 $\pm$ 0.033 & \underline{0.316 $\pm$ 0.013} & 5 \\
Spring-mass MPC & heldout & 2.085 $\pm$ 0.010 & 0.104 $\pm$ 0.065 & 0.920 $\pm$ 0.016 & 0.470 $\pm$ 0.014 & 5 \\
Two-mode spring-mass MPC & heldout & 2.009 $\pm$ 0.015 & 0.228 $\pm$ 0.101 & 0.961 $\pm$ 0.021 & 0.462 $\pm$ 0.013 & 5 \\
\textbf{Learned-Ctx MPC} & heldout & \underline{1.996 $\pm$ 0.009} & \underline{0.024 $\pm$ 0.017} & \underline{0.892 $\pm$ 0.019} & \textbf{0.176 $\pm$ 0.012} & 5 \\
\midrule
Pendulum MPC & heldout\_push & 2.270 $\pm$ 0.032 & \textbf{0.008 $\pm$ 0.011} & \textbf{0.812 $\pm$ 0.014} & 0.473 $\pm$ 0.015 & 5 \\
ZV input shaping & heldout\_push & \textbf{1.890 $\pm$ 0.000} & 0.972 $\pm$ 0.023 & 1.276 $\pm$ 0.026 & -- & 5 \\
Learned MPC (no context) & heldout\_push & 2.057 $\pm$ 0.008 & 0.184 $\pm$ 0.093 & 0.922 $\pm$ 0.026 & \underline{0.316 $\pm$ 0.013} & 5 \\
Spring-mass MPC & heldout\_push & 2.160 $\pm$ 0.049 & 0.132 $\pm$ 0.097 & 0.932 $\pm$ 0.018 & 0.470 $\pm$ 0.014 & 5 \\
Two-mode spring-mass MPC & heldout\_push & 2.029 $\pm$ 0.041 & 0.292 $\pm$ 0.125 & 0.972 $\pm$ 0.025 & 0.462 $\pm$ 0.013 & 5 \\
\textbf{Learned-Ctx MPC} & heldout\_push & \underline{2.024 $\pm$ 0.006} & \underline{0.120 $\pm$ 0.032} & \underline{0.915 $\pm$ 0.011} & \textbf{0.176 $\pm$ 0.012} & 5 \\
\midrule
Pendulum MPC & lin1 & 1.806 $\pm$ 0.005 & \underline{0.000 $\pm$ 0.000} & \textbf{0.898 $\pm$ 0.003} & 0.033 $\pm$ 0.000 & 5 \\
ZV input shaping & lin1 & \textbf{1.650 $\pm$ 0.000} & 0.000 $\pm$ 0.000 & 0.949 $\pm$ 0.000 & -- & 5 \\
Learned MPC (no context) & lin1 & 1.799 $\pm$ 0.009 & 0.004 $\pm$ 0.009 & 0.933 $\pm$ 0.010 & 0.066 $\pm$ 0.010 & 5 \\
Spring-mass MPC & lin1 & \underline{1.775 $\pm$ 0.006} & 0.012 $\pm$ 0.018 & 0.944 $\pm$ 0.005 & \textbf{0.032 $\pm$ 0.001} & 5 \\
Two-mode spring-mass MPC & lin1 & 1.777 $\pm$ 0.005 & 0.016 $\pm$ 0.009 & 0.944 $\pm$ 0.002 & \underline{0.032 $\pm$ 0.001} & 5 \\
\textbf{Learned-Ctx MPC} & lin1 & 1.797 $\pm$ 0.014 & \textbf{0.000 $\pm$ 0.000} & \underline{0.929 $\pm$ 0.007} & 0.046 $\pm$ 0.003 & 5 \\
\midrule
Pendulum MPC & push & 2.183 $\pm$ 0.014 & \textbf{0.000 $\pm$ 0.000} & \textbf{0.813 $\pm$ 0.013} & 0.135 $\pm$ 0.005 & 5 \\
ZV input shaping & push & \textbf{1.890 $\pm$ 0.000} & 0.988 $\pm$ 0.018 & 1.090 $\pm$ 0.008 & -- & 5 \\
Learned MPC (no context) & push & 2.002 $\pm$ 0.023 & 0.028 $\pm$ 0.033 & 0.918 $\pm$ 0.017 & \underline{0.092 $\pm$ 0.003} & 5 \\
Spring-mass MPC & push & 2.088 $\pm$ 0.032 & 0.072 $\pm$ 0.076 & 0.938 $\pm$ 0.013 & 0.136 $\pm$ 0.005 & 5 \\
Two-mode spring-mass MPC & push & \underline{1.981 $\pm$ 0.010} & 0.056 $\pm$ 0.026 & 0.922 $\pm$ 0.005 & 0.094 $\pm$ 0.007 & 5 \\
\textbf{Learned-Ctx MPC} & push & 2.012 $\pm$ 0.007 & \underline{0.024 $\pm$ 0.017} & \underline{0.903 $\pm$ 0.008} & \textbf{0.063 $\pm$ 0.004} & 5 \\
\bottomrule
\end{tabular}
}
\end{table*}

\begin{figure*}[t]\centering
\includegraphics[width=0.9\textwidth]{tradeoff.pdf}
\caption{Spill rate against mean transport time as each system's tuned scalar is swept (mean over 3 seeds). Left: nominal liquid, validation episodes used for tuning. Right: held-out liquids with a push, test episodes, three systems. Curves closer to the lower left are better.}\label{fig:tradeoff}
\end{figure*}

\begin{figure}[t]\centering
\includegraphics[width=\columnwidth]{trajectories.pdf}
\caption{One test episode on the nominal liquid (seed 0) for four systems at their tuned settings: commanded acceleration, cart velocity and the larger wall elevation.}\label{fig:traj}
\end{figure}

\subsection{H1: nominal liquid}
All systems are spill-free on the nominal test episodes at their tuned settings. The learned-context controller needs \rhval{main/learned-ctx-mpc/nominal/transport_time_s/mean:3}\,s and the pendulum controller \rhval{main/pendulum-mpc/nominal/transport_time_s/mean:3}\,s, a change of \rhval{cmp/main/pendulum-mpc/nominal/transport_time_s/rel_delta:pct1} relative to the pendulum (Welch $p=\,$\rhval{cmp/main/pendulum-mpc/nominal/transport_time_s/welch_p}); the per-seed ranges do not overlap. \textbf{H1 is supported.} Three qualifications matter more than the headline.

First, the size of the gap depends on the margin grid. The pendulum and the single-mode spring-mass model have almost the same prediction error (Table~\ref{tab:main}), yet the validation rule gives them margins of \rhval{select/tuned-setting/nominal/tuned_pendulum/mean:2} and \rhval{select/tuned-setting/nominal/tuned_springmass/mean:2}, because the pendulum controller spilled in one validation run at the larger value. Against the spring-mass controller (\rhval{main/spring-mass-mpc/nominal/transport_time_s/mean:3}\,s) the change is \rhval{cmp/main/spring-mass-mpc/nominal/transport_time_s/rel_delta:pct1}.

Second, the learned model is not the fastest model-based system. The two-mode spring-mass controller is spill-free at \rhval{main/two-mode-spring-mass-mpc/nominal/transport_time_s/mean:3}\,s, slightly faster than the learned controller (Welch $p=\,$\rhval{cmp/main/two-mode-spring-mass-mpc/nominal/transport_time_s/welch_p}). The learned model has the lower 20-step prediction error on the nominal liquid and runs at a larger margin (Table~\ref{tab:tuned}), but that did not translate into a shorter time.

Third, open-loop ZV input shaping is the fastest system on the nominal liquid, \rhval{main/zv-input-shaping/nominal/transport_time_s/mean:3}\,s, and it is also faster than the privileged oracle MPC (\rhval{reference/oracle-mpc-privileged/nominal/transport_time_s/mean:3}\,s, spill rate \rhval{reference/oracle-mpc-privileged/nominal/spill_rate/mean:pct1} at margin \rhval{select/tuned-setting/nominal/tuned_oracle/mean:2}). Two things explain this. The sampling optimiser is not exactly time-optimal: with a perfect model it still leaves time on the table. And shaping, being deterministic on a fixed liquid, can be tuned to a peak elevation of \rhval{main/zv-input-shaping/nominal/peak_disp/mean:3} freeboards with no validation spill. The model-related differences among MPC systems (a few hundredths of a second) are therefore of the same order as the loss due to the optimiser.

\subsection{H2: held-out liquids and pushes}
A push on the nominal liquid already separates the systems (task \emph{push} in Table~\ref{tab:main}): input shaping spills in \rhval{main/zv-input-shaping/push/spill_rate/mean:pct0} of episodes, the learned-context controller in \rhval{main/learned-ctx-mpc/push/spill_rate/mean:pct1}, the two-mode spring-mass controller in \rhval{main/two-mode-spring-mass-mpc/push/spill_rate/mean:pct1} and the pendulum controller in none.

The picture reverses for input shaping as soon as the liquid differs from the one it was tuned on: it spills in \rhval{main/zv-input-shaping/heldout/spill_rate/mean:pct0} of held-out episodes. Among MPC systems at their tuned margins on \emph{heldout\_push}, the learned-context controller takes \rhval{main/learned-ctx-mpc/heldout_push/transport_time_s/mean:3}\,s with a spill rate of \rhval{main/learned-ctx-mpc/heldout_push/spill_rate/mean:pct1}; the pendulum controller takes \rhval{main/pendulum-mpc/heldout_push/transport_time_s/mean:3}\,s with \rhval{main/pendulum-mpc/heldout_push/spill_rate/mean:pct1}. So at the tuned margins the pendulum controller spills \emph{less} (Welch $p=\,$\rhval{cmp/main/pendulum-mpc/heldout_push/spill_rate/welch_p}); it buys that with a longer transport: the difference in mean time, learned minus pendulum, is \rhval{cmp/main/pendulum-mpc/heldout_push/transport_time_s/delta:2}\,s. Margins tuned on the nominal liquid do not transfer for the learned model: a margin with no nominal validation spill gives a non-zero spill rate on held-out liquids even without a push (\rhval{main/learned-ctx-mpc/heldout/spill_rate/mean:pct1}).

H2 is stated at equal transport time, which the right panel of Fig.~\ref{fig:tradeoff} addresses. At a margin of 0.8 both controllers take about the same time (\rhval{derived/derived/nominal/hp_m080_learned_time/mean:3}\,s learned, \rhval{derived/derived/nominal/hp_m080_pendulum_time/mean:3}\,s pendulum) and the spill rates are \rhval{derived/derived/nominal/hp_m080_learned_spill/mean:pct1} and \rhval{derived/derived/nominal/hp_m080_pendulum_spill/mean:pct1}. Conversely, the pendulum controller reaches a spill rate of \rhval{derived/derived/nominal/hp_m070_pendulum_spill/mean:pct1} only at \rhval{derived/derived/nominal/hp_m070_pendulum_time/mean:3}\,s. \textbf{H2 is supported in the equal-time sense} (3 seeds), and not in the sense of a lower spill rate at the deployed settings. The context matters here: the no-context learned controller has a higher 20-step prediction error on held-out liquids (Table~\ref{tab:main}, $p=\,$\rhval{cmp/main/learned-mpc-no-context/heldout/pred_rmse_20/welch_p}) and its curve in Fig.~\ref{fig:tradeoff} lies between the other two; its spill-rate difference to the context model at the tuned margin is not significant on \emph{heldout\_push} ($p=\,$\rhval{cmp/main/learned-mpc-no-context/heldout_push/spill_rate/welch_p}).

The two-mode spring-mass controller, the fastest MPC system on the nominal liquid, is the least robust of them: \rhval{main/two-mode-spring-mass-mpc/heldout_push/spill_rate/mean:pct1} spills on \emph{heldout\_push}. Its extra accuracy is specific to the liquid it was identified on.

\subsection{H3: single linear mode}
With a single linear mode as ground truth the fitted models are structurally exact, and their 20-step error is lower than that of the learned model (Table~\ref{tab:main}). The learned-context and pendulum controllers take \rhval{main/learned-ctx-mpc/lin1/transport_time_s/mean:3}\,s and \rhval{main/pendulum-mpc/lin1/transport_time_s/mean:3}\,s (Welch $p=\,$\rhval{cmp/main/pendulum-mpc/lin1/transport_time_s/welch_p}), both without spills. \textbf{H3 is supported}: the advantage over the pendulum is within seed noise. Input shaping is again fastest (\rhval{main/zv-input-shaping/lin1/transport_time_s/mean:3}\,s); its tuned peak acceleration is the largest value of the grid we searched, so it might be faster still.
